\label{subsec:xi-certificate-coherence-orthogonal-axis-capacity-part3}

\begin{corollary}[Invisibility threshold under fixed resolution budget (bit/ledger calibration)]\label{cor:xi-offcritical-visibility-threshold-bit-budget}
Under the setting of Theorem \ref{thm:xi-offcritical-quadratic-radial-compression},
if one wishes to \emph{force} a positive defect contribution from this off-critical-line zero at some radius audit point \(\varrho\in(0,1)\) (i.e., require \(\varrho>|w_\rho|\), hence \(D_\zeta(\varrho)\ge \log(\varrho/|w_\rho|)>0\)),
then the endpoint resolution threshold must be satisfied
\[
1-\varrho\ \lesssim\ \frac{\delta}{\gamma^2}
\]
(the high-height scaling law of Corollary \ref{cor:app-horocycle-visible-resolution}).
In particular, if the resolution budget is parametrized by dyadic bit depth \(b\) such that
\(
1-\varrho\asymp 2^{-b}
\),
then the necessary condition can be equivalently written as
\begin{equation}\label{eq:xi-offcritical-bit-budget}
\boxed{\;
b\log 2\ \gtrsim\ 2\log|\gamma|+\log\frac{1}{\delta}+O(1)\; }.
\end{equation}
In other words: under fixed budget \(b\) there exists a height threshold such that when \(|\gamma|\) is sufficiently large, the second-order radial signal of the off-critical-line offset must fall into an indistinguishable scale;
in the pure phase regime (unitary slice), its semantic consequence stably manifests as "non-addressable" (\(\mathsf{NULL}\), Corollary \ref{cor:xi-null-second-order-radial-channel}).
\end{corollary}

\begin{proof}
By Corollary \ref{cor:app-horocycle-visible-resolution}, the necessary condition is \(1-\varrho\lesssim \delta/\gamma^2\).
Substituting \(1-\varrho\asymp 2^{-b}\) and taking logarithms yields
\(
b\log 2\gtrsim 2\log|\gamma|+\log(1/\delta)+O(1)
\).
The last sentence is given by "the second-order radial channel is erased by the pure phase regime" of Corollary \ref{cor:xi-null-second-order-radial-channel}.\qed
\end{proof}

The leading term \(2\log|\gamma|\) in equation \eqref{eq:xi-offcritical-bit-budget} embodies the irreducible demand for endpoint resolution from height growth;
this order is consistent with the "second-order radial channel" of Theorem \ref{thm:xi-offcritical-quadratic-radial-compression}.
On a finite height band \(|\gamma|\le T\), the same compression law also allows rewriting the existence of "off-slice" as a single-radius defect inequality certificate, thus obtaining a direct band exclusion conclusion.

\begin{theorem}[Finite sampling blind zone and band exclusion certificate: single-radius Jensen lower bound yields verifiable exclusion domain]\label{thm:xi-jensen-single-radius-band-exclusion}
Let \(T\ge 1\), \(\Delta\in(0,\tfrac12]\), and define
\[
  R(T,\Delta)^2:=\frac{T^2+(\Delta-1)^2}{T^2+(\Delta+1)^2}\in(0,1).
\]
Then for any off-critical-line zero \(\rho=\beta+\mathrm{i}\gamma\) satisfying \(|\gamma|\le T\), \(\delta=\tfrac12-\beta\ge\Delta\), its disk image \(w_\rho\) satisfies
\[
  |w_\rho|\le R(T,\Delta).
\]
Furthermore, for any audit radius \(\varrho\in(R(T,\Delta),1)\) there is a single-zero Jensen lower bound
\[
  D_\zeta(\varrho)\ \ge\ \log\frac{\varrho}{|w_\rho|}\ \ge\ \log\frac{\varrho}{R(T,\Delta)}
\]
(Proposition \ref{prop:app-jensen-single-zero-lower-bound}).
Therefore, if at some selected radius \(\varrho>R(T,\Delta)\) one verifies that
\[
  D_\zeta(\varrho)\ <\ \log\frac{\varrho}{R(T,\Delta)},
\]
this serves as an offline verifiable \emph{band exclusion certificate}: under this protocol calibration, there exists no off-critical-line zero satisfying \(|\gamma|\le T\) and \(\tfrac12-\beta\ge\Delta\).
\end{theorem}

\begin{proof}
By the closed form of Theorem \ref{thm:xi-offcritical-quadratic-radial-compression},
\[
|w_\rho|^2
:=1-\frac{4\delta}{\gamma^2+(1+\delta)^2}
:=\frac{\gamma^2+(1-\delta)^2}{\gamma^2+(1+\delta)^2}.
\]
Note that as \(\gamma^2\) increases, the right-hand side increases monotonically; as \(\delta\) increases, the right-hand side decreases monotonically (within the band \(\delta\in(0,1)\)).
Therefore under constraints \(|\gamma|\le T\), \(\delta\ge\Delta\),
\[
|w_\rho|^2
\le \frac{T^2+(1-\Delta)^2}{T^2+(1+\Delta)^2}
:=R(T,\Delta)^2,
\]
hence \(|w_\rho|\le R(T,\Delta)\).
When \(\varrho>|w_\rho|\), Proposition \ref{prop:app-jensen-single-zero-lower-bound} gives
\(
D_\zeta(\varrho)\ge \log(\varrho/|w_\rho|)
\); using \(|w_\rho|\le R(T,\Delta)\) yields the stated lower bound and certificate form.\qed
\end{proof}

\begin{corollary}[Explicit single-point certificate taking \texorpdfstring{$\varrho=\tfrac12(1+R(T,\Delta))$}{varrho=(1+R)/2}]\label{cor:xi-jensen-single-radius-band-exclusion-midradius}
Following the notation of Theorem \ref{thm:xi-jensen-single-radius-band-exclusion}, let
\(
\varrho:=\tfrac12\bigl(1+R(T,\Delta)\bigr)
\).
If one verifies that
\[
  D_\zeta(\varrho)\ <\ \log\frac{1+R(T,\Delta)}{2R(T,\Delta)},
\]
then it follows that in the band region \(|\gamma|\le T\) and \(\tfrac12-\beta\ge\Delta\) there exists no off-critical-line zero.
\end{corollary}

\begin{proof}
Substituting \(\varrho=\tfrac12(1+R(T,\Delta))\) into the certificate inequality of Theorem \ref{thm:xi-jensen-single-radius-band-exclusion} yields the result.\qed
\end{proof}

\paragraph{Falsifiable restatement of "constructing an off-critical-line zero": fourth register or failure witness}
In this paper's closure semantics, "constructing a zero not on \(\Re(s)=\tfrac12\)" is not equivalent to giving a numerical point;
it is equivalent to giving a \emph{verifiable address object} that can pass the closure checker (Definition \ref{def:multi_axis_address}), or giving a failure witness that such an address cannot exist within the protocol.

\begin{proposition}[Unique falsifiable interface for off-critical-line zero assertions: extended mode axis or output \texorpdfstring{$\mathsf{NULL}$}{NULL} failure witness]\label{prop:xi-offcritical-falsifiable-restatement}
Under unitary slice visible domain constraints and \(\mathbf{PW}\) type closure, any assertion claiming "there exists an off-critical-line zero (\(|u|\neq 1\) equivalently \(|w|<1\))" can only manifest within the system as one of the following two:
\begin{enumerate}
  \item \textbf{Mode axis extension (explicit radial record)}: Extend the readout/certificate type so that the mode axis \(a_{\mathrm{mode}}\) in the three-axis address prefix explicitly carries radial data (the \(\operatorname{proj}^{\mathrm e}\) calibration of Remark \ref{rem:xi-soundness-vs-completeness}); and at high height, this extension must pay the endpoint capacity lower bound given by Corollary \ref{cor:xi-offcritical-visibility-threshold-bit-budget}.
  \item \textbf{Failure witness (protocol lockdown)}: Output an offline verifiable failure witness proving that this assertion cannot be type-checked under the established visible domain/closure protocol, hence its readout must degrade to \(\mathsf{NULL}\) (Theorem \ref{thm:xi-offset-null-type-safety}; see also the trichotomy of \(\mathsf{NULL}\) in Proposition \ref{prop:xi-null-three-way}).
\end{enumerate}
Moreover: if one allows externalized side information with the goal of "reducing \(\mathsf{NULL}\)" to achieve injectivization, then this externalized axis must structurally be a residue/inverse limit axis (prime-register) rather than an additional continuous coordinate register decoupled from phase (Theorem \ref{thm:xi-pw-no-continuous-hair}).
\end{proposition}

\begin{proof}
If the mode axis is not extended, then \(|u|\neq 1\) (or \(|w|<1\)) falls outside the visible domain and is directly read out as \(\mathsf{NULL}\) by the filtering of Definition \ref{def:xi-visible-read-null}, with type-safe rejection given by Theorem \ref{thm:xi-offset-null-type-safety}.
If the mode axis is extended, then the extension-preserving construction of Remark \ref{rem:xi-soundness-vs-completeness} yields a non-degenerate extended readout; its necessary budget at high height is given by Corollary \ref{cor:xi-offcritical-visibility-threshold-bit-budget}.
Finally, "externalization must be residue/inverse limit axis rather than additional continuous register" is clause (2) of Theorem \ref{thm:xi-pw-no-continuous-hair}.\qed
\end{proof}

\begin{corollary}[Dichotomy of off-slice zero events: accept or \texorpdfstring{$\mathsf{NULL}$}{NULL}\,+\,offline witness]\label{cor:xi-offcritical-dichotomy-acceptable-or-null}
Under visible domain constraints locked by unitary slice and \(\mathbf{PW}\) closure, any attempt to objectify an "off-slice zero event" as a verifiable readout can only result in:
\begin{enumerate}
  \item \textbf{(Acceptance)} The event passes the type-checking verification interface of the unitary slice and is accepted, thus falling into the protocol-allowed visible type domain; in this sense, its zero parameter satisfies \(|u_L(s_0)|=1\);
  \item \textbf{(Falsifiable rejection)} The event has no address within the closure and is read out as \(\mathsf{NULL}\); and according to Theorem \ref{thm:xi-null-complete-trichotomy-offline}, this \(\mathsf{NULL}\) must be accompanied by an offline verifiable failure witness, whose support must be non-trivial in one of three orthogonal directions: "domain-external rejection/gluing torsion/hard resource insufficiency."
\end{enumerate}
Therefore there exists no third path "to circumvent \(\mathsf{NULL}\) by introducing an additional continuous coordinate register decoupled from phase": any continuous extension dimension that aims to reduce \(\mathsf{NULL}\) must, in the factor sense, be equivalent to externalizing the unique radial register \(\varrho(u)=\log|u|\) (Corollary \ref{cor:xi-unique-continuous-transverse-register}), otherwise it will still be compiled to \(\mathsf{NULL}\) by the closure type system.
\end{corollary}

\begin{proof}
The dichotomy comes from Proposition \ref{prop:xi-offcritical-falsifiable-restatement}: off-critical-line zero assertions can only appear within the system as "mode axis extension" or "\(\mathsf{NULL}\) failure witness."
When the assertion can be accepted within the pure phase regime, it must pass the visible domain filtering of the unitary slice (Definition \ref{def:xi-visible-read-null}), hence satisfying \(|u_L(s_0)|=1\); when the assertion cannot be type-checked, its readout is \(\mathsf{NULL}\), and the offline witness classification is obtained from Theorem \ref{thm:xi-null-complete-trichotomy-offline}.
The last sentence is given by the "no additional continuous hair" principle (Theorem \ref{thm:xi-pw-no-continuous-hair}) and the continuous axis uniqueness of Corollary \ref{cor:xi-unique-continuous-transverse-register}.\qed
\end{proof}

\paragraph{\texorpdfstring{$\tau$}{tau}-involution automorphism symmetry rigidity: \texorpdfstring{$\ZZ/2$}{Z/2} decomposition of Toeplitz-PSD, residual lower bound of Jensen defect, and localization of anomaly cohomology}
On the two-variable interface of \S\ref{subsec:xi-from-self-dual}, \(\tau:(z,u)\mapsto(uz,u^{-1})\) is an involution automorphism;
the self-dual condition \(\Delta=\Delta\circ\tau\) pulls back on any section of base \(L>1\) to the strict functional equation \(\Xi_{\Delta,L}(s)=\Xi_{\Delta,L}(1-s)\) (Proposition \ref{prop:xi-delta-L-functional-equation}).
Its fixed-point set corresponds to the unitary slice of the completion parameter: \(|u_L(s)|=1\iff \Re(s)=\tfrac12\) (Remark \ref{rem:xi-critical-line-unit-circle}).
Under this paper's visible domain filtering and "address precedes value" semantics, this fixed slice is promoted to an addressable regime: when \(X_{\mathrm{vis}}=\mathsf{Phase}\), points off the slice have no address at the type level, and their readout degrades to \(\mathsf{NULL}\) per protocol (Definition \ref{def:xi-visible-read-null}; Theorem \ref{thm:xi-offset-null-type-safety}), consistent with the calibration that recursive addressing only reselects events within the old-layer visible \(\sigma\)-algebra without expanding the visible domain (Proposition \ref{prop:recursive-addressing-refinement}; Corollary \ref{cor:recursive-addressing-visible-sigma-nonexpansion}).
\par
On the spectral side, Carathéodory positivity of the disk horizon logarithmic derivative is equivalent to the Toeplitz-PSD hierarchy \(\{\mathbf{T}_N\succeq 0\}_{N\ge 0}\) (Appendix Theorem \ref{thm:app-horizon-toeplitz-lmi}).
When the real structure \(c_n\in\RR\) is further locked, this hierarchy commutes with the reversal involution on finite truncations, hence splitting into even/odd two independent PSD sub-hierarchies (Proposition \ref{prop:app-horizon-toeplitz-z2-splitting}; Corollary \ref{cor:app-horizon-toeplitz-z2-cert-split});
thus any finite-level PSD failure evidence can be localized to some minimal \(N_\star\) and some sector sign \(\epsilon\in\{+,-\}\), giving a \(\ZZ/2\) classification of "failure modes."
\par
On the calculus side, the commutator failure residue of \(\mathbf{PW}\) is organized as an anomaly signature and promoted to a \(2\)-cohomology class (\S\ref{subsec:pom-anom-cohomology});
its visibility localization is given by Beck-Chevalley strictification: the anomaly is zero on visible twist channels (Proposition \ref{prop:pom-bc-quotient-strictify}), hence any nonzero anomaly coordinate can only be supported in directions non-trivial on \(N^{\mathrm{BC}}_m\) that are invisible (Corollary \ref{cor:pom-anom-supported-on-bc-kernel}).
\begin{remark}[Juxtaposed localization of Toeplitz sectors and anomaly class]\label{rem:xi-toeplitz-z2-anom-alignment}
Proposition \ref{prop:app-horizon-toeplitz-z2-splitting} gives the \(\ZZ/2\) sector decomposition of Toeplitz-PSD certificates under real structure lockdown;
Corollary \ref{cor:pom-anom-supported-on-bc-kernel} localizes the commutator failure of \(\mathbf{PW}\) as an anomaly cohomology class supported only on invisible role directions.
Their juxtaposition provides an implementation-independent obstruction coordinate system: if some assertion exhibits breaking in the unitary slice closure semantics, its spectral-side witness can be localized (at finite level) to a negative eigenvalue in some \(\epsilon\in\{+,-\}\) sector, while its calculus-side residue can be carried by nonzero coordinates of the anomaly class;
the two types of localization are not mutually substitutable, but can be recorded in parallel in the same certificate chain.
\par
In the multi-channel/operator-valued calibration, the above "channel-wise PSD + anomaly strictification" juxtaposition can be uniformly promoted to a completely positive horizon object valued in finite-dimensional \(C^\ast\)-algebra, whose block Toeplitz-PSD constraint is strictly equivalent to Fourier-Pontryagin (Peter-Weyl) decomposition; and strictification can be interpreted as the maximal quotient of "completely positive visible realization existence" (see \S\ref{subsubsec:xi-operator-valued-carath-toeplitz-complete-positivity}).
\end{remark}
In the disk completion interface, the so-called "leakage potential" can be taken as the relative quantity of logarithmic average potential, i.e., the Jensen defect (Definition \ref{def:app-jensen-defect}):
\[
D_F(\varrho)
:=\frac{1}{2\pi}\int_{0}^{2\pi}\log\bigl|F(\varrho e^{\mathrm{i}\theta})\bigr|\,\dd\theta
-\log|F(0)|
\qquad(0<\varrho<1).
\]
By the Jensen-Blaschke summation formula (Proposition \ref{prop:app-jensen-formula-defect}),
\[
D_F(\varrho)=\sum_{|a_k|<\varrho}\log\!\Bigl(\frac{\varrho}{|a_k|}\Bigr)\ \ge\ 0,
\]
therefore "any point-like zero appearing in the disk" will forcibly generate a strictly positive, quantifiable leakage lower bound;
the minimal leakage for a single zero is given by Proposition \ref{prop:app-jensen-single-zero-lower-bound}.
In this paper's projection readout semantics, this perfectly aligns with Proposition \ref{prop:xi-offcritical-falsifiable-restatement}: any point-like singularity deviating from the fixed slice \(|u|=1\) either appears as a positive contribution in the defect ledger \(D_\zeta(\varrho)\) form and triggers additional orthogonal certificate (residue/budget) externalization, or when refusing to externalize the mode axis is compiled by the protocol to a \(\mathsf{NULL}\) failure witness.
\par
Furthermore, the "fixed slice" is not only the fixed set of automorphisms, but also the unique equilibrium surface of logarithmic average potential.
The following proposition gives a line-by-line verifiable inversion symmetry law, which directly translates the self-dual property of "inner and outer domains as mirror images" into the evenness rigidity of average potential.

\begin{proposition}[Logarithmic average potential inversion law under self-dual symmetry]\label{prop:xi-selfdual-mean-potential-inversion}
Let \(A=\{u\in\CC:\ r_0<|u|<r_0^{-1}\}\) and \(G\) be analytic on \(A\) with \(G\not\equiv 0\).
Suppose there exists an integer \(m\) such that for all \(u\in A\),
\[
G(u)=u^{m}\,\overline{G(1/\overline{u})}.
\]
For \(r\in(r_0,r_0^{-1})\) define the logarithmic average potential
\[
\mathcal{M}_G(r):=\frac{1}{2\pi}\int_{0}^{2\pi}\log\bigl|G(re^{\mathrm{i}\theta})\bigr|\,\dd\theta.
\]
Then for all \(r\in(r_0,r_0^{-1})\) the strict identity holds
\[
\boxed{\ \mathcal{M}_G(r)=m\log r+\mathcal{M}_G(1/r)\ }.
\]
Equivalently, letting \(r=e^t\), then \(t\mapsto \mathcal{M}_G(e^t)-\frac{m}{2}t\) is an even function, with the unique possible equilibrium slice being \(t=0\), i.e., \(|u|=1\).
\end{proposition}

\begin{proof}
Take \(u=re^{\mathrm{i}\theta}\), then \(1/\overline{u}=r^{-1}e^{\mathrm{i}\theta}\).
By assumption,
\(
|G(re^{\mathrm{i}\theta})|
:=|u|^{m}\,|G(r^{-1}e^{\mathrm{i}\theta})|
:=r^{m}|G(r^{-1}e^{\mathrm{i}\theta})|
\).
Taking logarithms on both sides and averaging over \(\theta\) yields the stated identity; substituting \(r=e^t\) and rearranging yields the evenness and equilibrium slice conclusion.\qed
\end{proof}

The above inversion law gives the strict version of "\(\Re(s)=\tfrac12\) is the symmetry axis": the symmetry is not interchange of coordinate axes, but inner and outer domains being mirror images under self-dual involution; hence its fixed set \(|u|=1\) becomes the unique equilibrium surface of (normalized) logarithmic average potential.
\par
Building on this, synthesizing Theorem \ref{thm:xi-orthogonal-slice-structure}, Theorem \ref{thm:xi-offcritical-quadratic-radial-compression} and its corollary chains, the meaning of "unitary slice (fixed visible slice of self-dual projection)" under protocol calibration can be condensed into the following four conclusive propositions:
\begin{enumerate}
  \item \textbf{(Endpoint payload decomposition)}
  The off-critical-line offset \(\delta\) and height \(\gamma\) appear as additive loads in the endpoint layer depth (Corollary \ref{cor:xi-offcritical-endpoint-depth-log-decomposition}), unified by the precision potential \(\mathcal{P}\) as ledger coordinate.
  \item \textbf{(Complexity barrier)}
  Moving from the phase axis of the unitary slice into the off-critical-line direction, the endpoint layer bit budget jumps from \(\Theta(\log T)\) to \(\Theta(2\log T)\), subject to explicit lower bound constraints (Theorem \ref{thm:xi-offcritical-endpoint-resolution-lower-bound} and Remark \ref{rem:xi-unitary-vs-offcritical-bit-gap}).
  \item \textbf{(Inner factor visibility)}
  Under Hardy/Smirnov condition calibration, off-critical-line zeros are equivalent to Blaschke inner factor data; this data cannot be generated from boundary modulus within the pure phase visible interface, hence its verification semantics must forcibly occupy an orthogonal register (Corollary \ref{cor:xi-inner-factor-orthogonal-register}).
  \item \textbf{(Endpoint phase certificate)}
  Under Hardy/Smirnov condition calibration, the inner factor can define a phase current measure \(\mu_{\mathrm{in}}\), whose endpoint density \(j_{\mathrm{in}}(\pi)\) gives a non-cancellable positive lower bound for in-disk zeros; and for the total off-critical-line offset satisfies \(j_{\mathrm{in}}(\pi)\ge \sum_k\delta_k\) (Theorem \ref{thm:app-inner-phase-current-poisson} and Corollary \ref{cor:app-endpoint-phase-current-offcritical}).
  Under the condition of no singular inner factor \(S\equiv 1\), the endpoint phase current zero certificate \(j_{\mathrm{in}}(\pi)=0\) forces \(B\equiv 1\), hence deducing in-disk zero exclusion and returning to RH (Corollary \ref{cor:app-endpoint-phase-certificate}).
\end{enumerate}
\par
\paragraph{Endpoint absorption coefficient: single-point flux compression of in-disk zero set}
The above endpoint phase current density \(j_{\mathrm{in}}(\pi)\) (in the finite Blaschke case, the sum of endpoint Poisson weights) not only gives a "non-cancellable positive lower bound," but also provides a single-point readable radial decay rate certificate: the in-disk zero set can be compressed into the absorption coefficient at endpoint \(-1\), read out by the first-order decay rate as \(w=-r\uparrow-1\).

\begin{proposition}[Endpoint absorption coefficient and radial absorption law]\label{prop:xi-endpoint-absorption-coefficient}
Let \(B\) be a finite Blaschke product with zero set \(\{a_k\}_{k=1}^{N}\subset\mathbb{D}\) (counted with multiplicity).
Define the absorption coefficient at endpoint \(-1\)
\[
\Phi_{-1}(B):=\sum_{k=1}^{N}P_{a_k}(-1)=\sum_{k=1}^{N}\frac{1-|a_k|^2}{|1+a_k|^2}\ \in[0,\infty).
\]
Then as \(r\uparrow 1\), along the endpoint radial \(w=-r\) the asymptotic law holds
\[
\log|B(-r)|=-(1-r)\,\Phi_{-1}(B)+o(1-r).
\]
In particular, \(\Phi_{-1}(B)=0\) if and only if \(B\equiv 1\).
\end{proposition}

\begin{proof}
See Appendix Proposition \ref{prop:app-endpoint-blaschke-radial-absorption}.\qed
\end{proof}

\begin{theorem}[Endpoint Julia-Carathéodory indicator equals absorption coefficient: angular derivative form of \texorpdfstring{$\Phi_{-1}$}{Phi-(-1)}]\label{thm:xi-endpoint-julia-indicator-equals-absorption}
Under the setting of Proposition \ref{prop:xi-endpoint-absorption-coefficient}, define the Julia-Carathéodory indicator at endpoint \(-1\)
\[
\mathcal{J}_{-1}(B):=\lim_{r\uparrow 1}\frac{1-|B(-r)|^2}{1-r^2}.
\]
Then this limit exists and satisfies
\[
\boxed{\;\mathcal{J}_{-1}(B)=\Phi_{-1}(B)\;}
\]
and (the angular derivative exists and is finite)
\[
\boxed{\;\Phi_{-1}(B)=\Re\!\left(-\frac{B'(-1)}{B(-1)}\right).\;}
\]
\end{theorem}

\begin{proof}
By the radial absorption law of Proposition \ref{prop:xi-endpoint-absorption-coefficient},
\(
\log|B(-r)|=-(1-r)\Phi_{-1}(B)+o(1-r)
\).
Hence
\[
|B(-r)|^2=\exp\!\bigl(2\log|B(-r)|\bigr)
=\exp\!\bigl(-2(1-r)\Phi_{-1}(B)+o(1-r)\bigr)
=1-2(1-r)\Phi_{-1}(B)+o(1-r),
\]
thus
\(
1-|B(-r)|^2=2(1-r)\Phi_{-1}(B)+o(1-r)
\).
Note \(1-r^2=2(1-r)+o(1-r)\), hence \(\mathcal{J}_{-1}(B)=\Phi_{-1}(B)\).
\par
On the other hand, the finite Blaschke product is analytic in a closed disk neighborhood and \(|B(-1)|=1\), so
\(
B(-r)=B(-1)+(1-r)B'(-1)+o(1-r)
\);
substituting and expanding gives
\(
1-|B(-r)|^2=-2(1-r)\Re\!\bigl(B'(-1)/B(-1)\bigr)+o(1-r)
\),
dividing by \(1-r^2\) and letting \(r\uparrow 1\) yields
\(
\mathcal{J}_{-1}(B)=\Re\!\bigl(-B'(-1)/B(-1)\bigr)
\).
Combining with the first step yields the conclusion.\qed
\end{proof}

\begin{corollary}[Additive conservation and single-point tomography of endpoint absorption coefficient]\label{cor:xi-endpoint-absorption-additive-tomography}
Under the setting of Proposition \ref{prop:xi-endpoint-absorption-coefficient}, the following identities and verifiable readouts hold:
\begin{enumerate}
  \item \textbf{(Additive conservation of endpoint Poisson-Busemann mass)}
  \[
  \Phi_{-1}(B)=\sum_{k=1}^{N}\mathcal{B}_{-1}(a_k)
  \qquad(\text{counted with multiplicity}),
  \]
  where \(\mathcal{B}_{-1}(w)=\frac{1-|w|^2}{|1+w|^2}\) is the exponential Busemann coordinate at endpoint \(-1\) (Proposition \ref{prop:app-busemann-poisson-minus1}).
  \item \textbf{(Single-point tomography: reading out total offset along one radial)}
  \[
  \Phi_{-1}(B)=\lim_{r\uparrow 1}\frac{-\log|B(-r)|}{1-r}.
  \]
  \item \textbf{(Certificating total off-slice offset)}
  If further assuming \(a_k=w_{\rho_k}\) are given by disk images of off-slice zeros \(\rho_k\),
  and letting \(\delta_k:=\tfrac12-\Re(\rho_k)=\mathcal{B}_{-1}(w_{\rho_k})\) (Theorem \ref{thm:xi-offcritical-quadratic-radial-compression}),
  then
  \[
  \Phi_{-1}(B)=\sum_{k=1}^{N}\delta_k.
  \]
\end{enumerate}
\end{corollary}

\begin{proof}
Clause (1) follows directly from \(P_{a}(-1)=\frac{1-|a|^2}{|1+a|^2}=\mathcal{B}_{-1}(a)\) (Proposition \ref{prop:app-busemann-poisson-minus1}) and the definition \(\Phi_{-1}(B)=\sum_k P_{a_k}(-1)\).
Clause (2) follows by dividing both sides of the asymptotic formula in Proposition \ref{prop:xi-endpoint-absorption-coefficient} by \(1-r\) and letting \(r\uparrow 1\).
Clause (3) follows by combining clause (1) with \(\delta_k=\mathcal{B}_{-1}(w_{\rho_k})\) (Theorem \ref{thm:xi-offcritical-quadratic-radial-compression}).\qed
\end{proof}

\begin{theorem}[Completeness of "true resonance exclusion" by endpoint single-point certificate]\label{thm:xi-endpoint-single-point-certificate-complete}
Under Hardy/Smirnov calibration with no singular inner factor \(S\equiv 1\), suppose in-disk point-like data is carried by finite Blaschke factor \(B\).
Then the following statements are equivalent:
\begin{enumerate}
  \item \(B\equiv 1\) (no in-disk zeros/no true resonances);
  \item Endpoint phase current density \(j_{\mathrm{in}}(\pi)=0\) (Corollary \ref{cor:app-endpoint-phase-certificate});
  \item Endpoint absorption coefficient \(\Phi_{-1}(B)=0\) (Proposition \ref{prop:xi-endpoint-absorption-coefficient});
  \item Endpoint Julia-Carathéodory indicator \(\mathcal{J}_{-1}(B)=0\) (Theorem \ref{thm:xi-endpoint-julia-indicator-equals-absorption}).
\end{enumerate}
And in the finite Blaschke case there is strict identity
\[
\boxed{\;
j_{\mathrm{in}}(\pi)=\Phi_{-1}(B)=\sum_k\mathcal{B}_{-1}(a_k)\; }
\]
(see Theorem \ref{thm:app-inner-phase-current-poisson} and Corollary \ref{cor:xi-endpoint-absorption-additive-tomography}).
Therefore, the existence of "true resonances" (in-disk zeros/stable domain poles) can be compressed into a single-point non-negative scalar certificate at endpoint \(-1\): any one of the above quantities is strictly positive if and only if \(B\not\equiv 1\).
\end{theorem}

\begin{proof}
\((1)\Leftrightarrow(3)\) by the last sentence of Proposition \ref{prop:xi-endpoint-absorption-coefficient};
\((3)\Leftrightarrow(4)\) by the equation \(\mathcal{J}_{-1}(B)=\Phi_{-1}(B)\) of Theorem \ref{thm:xi-endpoint-julia-indicator-equals-absorption};
\((2)\Leftrightarrow(1)\) by Corollary \ref{cor:app-endpoint-phase-certificate} (under \(S\equiv 1\) the endpoint phase current zero certificate forces \(B\equiv 1\));
the final identity follows by combining the endpoint phase current density equals endpoint Poisson weight sum for finite Blaschke case (Theorem \ref{thm:app-inner-phase-current-poisson}) with \(\Phi_{-1}(B)=\sum_k\mathcal{B}_{-1}(a_k)\) from Corollary \ref{cor:xi-endpoint-absorption-additive-tomography}.\qed
\end{proof}

\paragraph{Endpoint Tomography Principle: An Auditable Inversion Chain Recovering Full ``Off-Axis Zero'' Data from a Single Radial Modulus Profile}\label{par:xi-endpoint-tomography-principle}
Under the Hardy/Smirnov regime, the disk-completed function \(F_{\zeta}\) admits an inner--outer factorization \(F_{\zeta}=BSO\) (Appendix \S\ref{subsec:unit-disk-inner-outer}).
The Blaschke factor \(B\) encodes, via its disk-interior zero set \(\{a_k\}\subset\mathbb{D}\), all discrete degrees of freedom corresponding to ``off-critical-line deviations''; its absorption coefficient at the endpoint \(-1\),
\(\Phi_{-1}(B)=\sum_k P_{a_k}(-1)\),
is read off from the first-order decay rate along a single radial direction (Proposition \ref{prop:xi-endpoint-absorption-coefficient}; see also Corollary \ref{cor:xi-endpoint-absorption-additive-tomography}).
In the case where the singular inner factor is trivial \(S\equiv 1\) and \(B\) is a finite Blaschke product, this quantity also equals the endpoint phase current density \(j_{\mathrm{in}}(\pi)\) (see \eqref{eq:app-endpoint-phase-current-lowerbound}).
This paragraph establishes a stronger tomographic result: under the finite-defect (finite Blaschke) regime, the \emph{entire endpoint radial modulus profile}
\(
r\mapsto |B(-r)|
\)
not only yields the total \(\sum_k\delta_k\), but also carries the complete data \((\gamma_k,\delta_k)\), and possesses inversion uniqueness in the sense of analytic continuation.

\begin{definition}[Endpoint Radial Absorption Profile]\label{def:xi-endpoint-radial-absorption-profile}
Let \(B\) be a finite Blaschke product.
Define the \textbf{radial absorption profile} at the endpoint \(-1\) as
\begin{equation}\label{eq:xi-endpoint-radial-absorption-profile}
\boxed{\;
\mathcal{A}_B(r):=-\frac{\dd}{\dd r}\log|B(-r)|
\qquad(0<r<1).\;}
\end{equation}
\end{definition}

\begin{theorem}[Full Kernel Representation of the Endpoint Radial Absorption: A One-Dimensional Profile Carries Two-Dimensional Zero Data]\label{thm:xi-endpoint-absorption-kernel-representation}
Let \(B\) be a finite Blaschke product with zero multiset \(\{a_k\}_{k=1}^{N}\subset\mathbb{D}\) (counted with multiplicity).
Let
\[
t_k:=\mathcal{C}^{-1}(a_k)=\gamma_k+\mathrm{i}\delta_k\in\mathbb{H},
\qquad (\gamma_k\in\RR,\ \delta_k>0),
\]
where \(\mathcal{C}\) is the fixed Cayley compactification of this paper (from \(\mathbb{H}\) to \(\mathbb{D}\); see Proposition \ref{prop:app-busemann-poisson-minus1} and Corollary \ref{cor:app-offcritical-horocycle-busemann}).
Then for any \(r\in(0,1)\) there is an explicit kernel decomposition
\begin{equation}\label{eq:xi-endpoint-absorption-kernel-sum}
\boxed{\;
\mathcal{A}_B(r)
=\sum_{k=1}^{N}
\frac{4\delta_k}{(1-r)^2\gamma_k^2+\bigl(1+r+(1-r)\delta_k\bigr)^2}\; }.
\end{equation}
Furthermore, the endpoint limit satisfies
\begin{equation}\label{eq:xi-endpoint-absorption-kernel-limit}
\boxed{\;
\lim_{r\uparrow 1}\mathcal{A}_B(r)
=\sum_{k=1}^{N}\delta_k
=\sum_{k=1}^{N}P_{a_k}(-1)
=\Phi_{-1}(B)\; }.
\end{equation}
\end{theorem}

\begin{proof}
Write \(B\) as a product of Blaschke factors \(B=\prod_{k=1}^{N}b_{a_k}\) (ignoring the constant phase factor).
For a single factor \(b_a(w)=\frac{a-w}{1-\overline{a}\,w}\) evaluated along the real radial direction \(w=-r\), one has
\[
\frac{\dd}{\dd r}\log|b_a(-r)|
=\Re\!\left(\frac{1}{a+r}-\frac{\overline a}{1+r\overline a}\right),
\]
so that \(\mathcal{A}_B(r)=-\sum_k \frac{\dd}{\dd r}\log|b_{a_k}(-r)|\).
\par
Let \(a=\mathcal{C}(t)\) with \(t=\gamma+\mathrm{i}\delta\in\mathbb{H}\).
Using the closed-form expression for \(\mathcal{C}\) (equivalent to the Cayley choice where \(\mathcal{B}_{-1}(w)=\Im(\mathcal{C}^{-1}(w))\) in Proposition \ref{prop:app-busemann-poisson-minus1}), one can combine the denominators and simplify to obtain
\[
-\frac{\dd}{\dd r}\log|b_a(-r)|
=\frac{4\delta}{(1-r)^2\gamma^2+\bigl(1+r+(1-r)\delta\bigr)^2}.
\]
Summing over \(k=1,\dots,N\) yields \eqref{eq:xi-endpoint-absorption-kernel-sum}.
\par
Taking \(r\uparrow 1\), the denominator tends to \(4\), so each term tends to \(\delta_k\), whence
\(\lim_{r\uparrow 1}\mathcal{A}_B(r)=\sum_k\delta_k\).
By \(P_{a_k}(-1)=\mathcal{B}_{-1}(a_k)=\Im(\mathcal{C}^{-1}(a_k))=\delta_k\) (\eqref{eq:app-poisson-kernel-endpoint-minus1}) and the definition \(\Phi_{-1}(B)=\sum_k P_{a_k}(-1)\) (Proposition \ref{prop:xi-endpoint-absorption-coefficient}), one obtains \eqref{eq:xi-endpoint-absorption-kernel-limit}. \qed
\end{proof}

\begin{theorem}[Endpoint Jet Invariants: Low-Order Derivatives of \texorpdfstring{$\mathcal{A}_B$}{A\_B} at \texorpdfstring{$r=1$}{r=1} Yield Verifiable ``Moments'']\label{thm:xi-endpoint-absorption-jet-moments}
Under the finite Blaschke setting of Theorem \ref{thm:xi-endpoint-absorption-kernel-representation}, the function \(\mathcal{A}_B\) is smooth in a left neighborhood of \(r\uparrow 1\) and satisfies the endpoint jet identities
\[
\boxed{\;
\mathcal{A}_B(1)=\sum_{k}\delta_k,\qquad
\mathcal{A}_B'(1)=\sum_{k}\delta_k(\delta_k-1),\qquad
\mathcal{A}_B''(1)=\frac12\sum_{k}\delta_k\bigl(3(\delta_k-1)^2-\gamma_k^2\bigr).\;}
\]
When \(\delta_k\in(0,\tfrac12)\) (off-critical-line deviation regime), one has \(\mathcal{A}_B'(1)<0\); moreover, the second-order term
\(-\tfrac12\sum_k\delta_k\gamma_k^2\)
provides a second-order suppression response for large heights \(|\gamma|\).
\end{theorem}

\begin{proof}
By \eqref{eq:xi-endpoint-absorption-kernel-sum}, \(\mathcal{A}_B\) is a finite sum of rational functions and hence smooth in a one-sided neighborhood of \(r=1\).
For the single-term kernel
\(
K_{(\gamma,\delta)}(r)=\frac{4\delta}{(1-r)^2\gamma^2+(1+r+(1-r)\delta)^2}
\),
expanding in \(u=1-r\) yields
\[
K_{(\gamma,\delta)}(1)=\delta,\quad
K_{(\gamma,\delta)}'(1)=\delta(\delta-1),\quad
K_{(\gamma,\delta)}''(1)=\tfrac12\,\delta\bigl(3(\delta-1)^2-\gamma^2\bigr),
\]
and summing over \(k\) gives the result. \qed
\end{proof}

\begin{theorem}[Inversion Uniqueness in the Finite-Defect Regime: \texorpdfstring{$\mathcal{A}_B$}{A\_B} Single-Variable Data Determines the Zero Multiset]\label{thm:xi-endpoint-absorption-inversion-uniqueness}
Under the setting of Theorem \ref{thm:xi-endpoint-absorption-kernel-representation}, let
\[
q:=\frac{1-r}{1+r}\in(0,1),
\qquad r=\frac{1-q}{1+q}.
\]
Then \(q\mapsto \mathcal{A}_B\!\bigl(\frac{1-q}{1+q}\bigr)\) is a rational function in the complex \(q\)-plane, whose pole set is completely determined by \(\{t_k=\gamma_k+\mathrm{i}\delta_k\}\): for each \(k\) there is a conjugate pair of poles
\begin{equation}\label{eq:xi-endpoint-absorption-q-poles}
\boxed{\;
q_{k,\pm}=\frac{-\delta_k\pm \mathrm{i}\gamma_k}{\delta_k^2+\gamma_k^2}
=-\frac{\delta_k\mp \mathrm{i}\gamma_k}{\delta_k^2+\gamma_k^2}\;}
\end{equation}
with residues depending only on \(\delta_k\).
Therefore:
\begin{enumerate}
  \item If two finite Blaschke products \(B_1,B_2\) satisfy \(\mathcal{A}_{B_1}(r)=\mathcal{A}_{B_2}(r)\) on some nontrivial interval, then their zero multisets in Cayley coordinates must agree (counted with multiplicity):
  \[
  \{\mathcal{C}^{-1}(a_k^{(1)})\}=\{\mathcal{C}^{-1}(a_k^{(2)})\}.
  \]
  \item In particular, the endpoint radial modulus profile \(r\mapsto |B(-r)|\) (equivalently, \(\mathcal{A}_B\)) determines all parameters \((\gamma_k,\delta_k)\) in the finite-defect regime, up to the natural indistinguishability arising from conjugate symmetry and permutations; this indistinguishability is entirely compatible with conjugate pairing and self-dual symmetry of the zeros.
\end{enumerate}
\end{theorem}

\begin{proof}
Substituting \(r=\frac{1-q}{1+q}\) into \eqref{eq:xi-endpoint-absorption-kernel-sum} and using the identities
\(
1-r=\frac{2q}{1+q},\ 1+r=\frac{2}{1+q}
\),
one obtains
\[
\mathcal{A}_B\!\Bigl(\frac{1-q}{1+q}\Bigr)
=\sum_{k=1}^{N}
\frac{\delta_k(1+q)^2}{q^2\gamma_k^2+(1+q\delta_k)^2},
\]
which is a rational function. Its poles are given by the vanishing of the denominator:
\(
q^2\gamma_k^2+(1+q\delta_k)^2=0
\)
is equivalent to
\(
(\delta_k^2+\gamma_k^2)q^2+2\delta_k q+1=0
\),
solving which yields \eqref{eq:xi-endpoint-absorption-q-poles}, with the conjugate pair following from conjugate symmetry.
The dependence of the residues on \(\delta_k\) alone is verified by direct substitution.
\par
If \(\mathcal{A}_{B_1}=\mathcal{A}_{B_2}\) on an interval, then the corresponding rational functions in the \(q\)-variable agree on an open set and hence, by analytic continuation (the identity theorem), agree everywhere as meromorphic functions; consequently their pole sets and multiplicities must coincide, so \(\{(\gamma_k,\delta_k)\}\) (equivalently \(\{t_k\}\)) agree with multiplicity, establishing conclusion (1).
Conclusion (2) follows from \(\mathcal{A}_B=-\frac{\dd}{\dd r}\log|B(-r)|\) and integration to recover \(|B(-r)|\) (combined with the normalization \(B(0)\neq 0\)). \qed
\end{proof}

\begin{corollary}[C-Finite Rigidity of Endpoint Profile Coefficients: Constant-Coefficient Recurrence and Hankel Rank Certificate]\label{cor:xi-endpoint-profile-cfinite-hankel-rank}
Under the setting of Theorem \ref{thm:xi-endpoint-absorption-inversion-uniqueness}, let
\[
F(q):=\mathcal{A}_B\!\Bigl(\frac{1-q}{1+q}\Bigr)=\sum_{n\ge 0}c_n q^n\qquad(|q|\ \text{sufficiently small}).
\]
Then \(F\) is a rational function with denominator given by
\[
\boxed{\ D(q):=\prod_{k=1}^{N}\bigl((\delta_k^2+\gamma_k^2)q^2+2\delta_k q+1\bigr)\in\ZZ[q]\ }.
\]
Therefore the coefficient sequence \(\{c_n\}_{n\ge 0}\) is a strict C-finite sequence: there exists a linear recurrence of order \(2N\) with constant coefficients such that for all sufficiently large \(n\),
\[
c_{n+2N}=\sum_{j=0}^{2N-1}\alpha_j\,c_{n+j}\qquad(\alpha_j\in\QQ).
\]
If we further assume that the pole pairs \(\{q_{k,\pm}\}\) are pairwise distinct (equivalent to the condition that \(\{t_k=\gamma_k+\mathrm{i}\delta_k\}\) are distinct in \(\mathbb{H}\) with multiplicity), then the above order is minimal; equivalently, the Hankel matrix
\(
\mathsf{H}=(c_{i+j})_{0\le i,j\le 2N}
\)
has rank exactly \(2N\), so that \(N\) can be read off precisely via finite-truncation linear algebra certificates.
\end{corollary}
\begin{proof}
By the rational expression in Theorem \ref{thm:xi-endpoint-absorption-inversion-uniqueness}, the poles of \(F\) are precisely the roots of the quadratic factors
\(
(\delta_k^2+\gamma_k^2)q^2+2\delta_k q+1
\),
namely \(q_{k,\pm}\). Taking \(D(q)\) to be the product of these quadratic factors ensures that \(D(q)F(q)\) is a polynomial, so \(F\) is a rational function and its Taylor coefficients satisfy a constant-coefficient linear recurrence given by the coefficients of \(D\) (obtained by comparing coefficients term-by-term in the power series \(D(q)F(q)\)).
\par
When the poles are pairwise distinct and no factors cancel, \(D\) is the minimal denominator (up to scalar), whose degree \(2N\) is the minimal recurrence order; and the minimal recurrence order is equivalent to the Hankel rank characterization. \qed
\end{proof}

\begin{corollary}[Prony Invertibility of Endpoint Data: \texorpdfstring{$4N$}{4N} Coefficients Recover All \texorpdfstring{$(\gamma_k,\delta_k)$}{(gamma,delta)}]\label{cor:xi-endpoint-profile-prony-invertibility}
Under the setting of Theorem \ref{thm:xi-endpoint-absorption-inversion-uniqueness}, if we further assume that the pole pairs \(\{q_{k,\pm}\}\) are pairwise distinct, then given any \(4N\) consecutive Taylor coefficients \(c_0,c_1,\dots,c_{4N-1}\) of \(F(q)\) at \(q=0\), one can in principle uniquely recover the parameter multiset \(\{(\gamma_k,\delta_k)\}_{k=1}^{N}\) (up to the natural indistinguishability due to conjugate pairing and permutations).
\end{corollary}
\begin{proof}
By Corollary \ref{cor:xi-endpoint-profile-cfinite-hankel-rank}, the minimal denominator \(D(q)\) of \(F\) has degree \(2N\), so \(\{c_n\}\) satisfies a minimal recurrence of order \(2N\); this recurrence (equivalently, the annihilating polynomial) can be uniquely recovered from \(c_0,\dots,c_{4N-1}\) via Hankel linear systems/Prony null-space methods (or equivalently, the Berlekamp--Massey algorithm), thereby yielding \(D(q)\) (normalized to have leading coefficient \(1\)).
Factoring \(D\) into quadratic factors recovers all poles \(q_{k,\pm}\); then by the sum-product relations in \eqref{eq:xi-endpoint-absorption-q-poles} one can solve for \((\gamma_k,\delta_k)\).
Finally, the residues (equivalently, the individual term coefficients) can be uniquely determined by writing \(F\) in partial fraction form and matching the first \(2N\) coefficients via a linear system. \qed
\end{proof}

\begin{corollary}[Endpoint Certificate Chain for \texorpdfstring{$F_\zeta$}{F\_zeta}: From a Single Radial Profile to Off-Line Zero Existence]\label{cor:xi-endpoint-tomography-certificate-chain}
Under the factorization regime of Appendix \S\ref{subsec:unit-disk-inner-outer}, if one further assumes \(S\equiv 1\) (no singular inner factor), then the endpoint radial profile \(r\mapsto |B(-r)|\) is jointly determined by the endpoint radial modulus profile of \(F_\zeta\) and the outer factor profile:
\[
\log|F_\zeta(-r)|=\log|B(-r)|+\log|O(-r)|.
\]
Therefore, under the condition that the radial log-derivative of the outer factor \(O\) can be auditably controlled and stripped off, non-vanishing of \(\mathcal{A}_B\) provides an endpoint certificate for the existence of off-line zeros; moreover, in the finite-defect approximation (finite Blaschke), \(\mathcal{A}_B\) can even invert all \((\gamma_k,\delta_k)\), thereby transcribing ``off-line zeros'' from two-dimensional hidden data into an auditable spectral fingerprint of a single-variable profile function (Theorem \ref{thm:xi-endpoint-absorption-inversion-uniqueness}).
\end{corollary}

\endinput


\paragraph{Height compression and endpoint decoupling: visible scale splitting of two certificate channels}
Theorem \ref{thm:xi-offcritical-quadratic-radial-compression} shows that for the disk image \(w_\rho\) of off-slice zero \(\rho=\beta+\mathrm{i}\gamma\), the radial gap satisfies
\(
1-|w_\rho|^2=\frac{4\delta}{\gamma^2+(1+\delta)^2}
\)
and compresses quadratically with \(|\gamma|\);
therefore if one attempts to read out the contribution of this zero at a fixed radius using only the radial gap channel (Jensen circular average), one must enter the endpoint boundary layer resolution \(1-\varrho\lesssim \delta/\gamma^2\), thus triggering an unavoidable resolution/bit budget threshold (Corollary \ref{cor:xi-offcritical-visibility-threshold-bit-budget}; see also the barrier form of Theorem \ref{thm:xi-jensen-quadratic-resolution-barrier}).
\par
In contrast, the offset coordinate \(\delta=\mathcal{B}_{-1}(w_\rho)\) of the same zero is the horocycle layer label at endpoint \(-1\) (Proposition \ref{prop:app-busemann-poisson-minus1} and Theorem \ref{thm:xi-offcritical-quadratic-radial-compression}), which does not decay with \(|\gamma|\);
under the calibration "offset only carried by Blaschke data," it directly enters the additive conservation quantity \(\Phi_{-1}(B)=\sum_k \delta_k\) of the endpoint absorption coefficient and can be read out by single-point radial decay rate (Corollary \ref{cor:xi-endpoint-absorption-additive-tomography}).
Therefore, the verifiable information of "off-slice offset" naturally splits into two orthogonal certificate channels in the closed-loop interface: one radial gap channel that undergoes quadratic radial compression with height, and another endpoint channel that decouples from height, carried by endpoint Poisson-Busemann mass and single-point flux.

\begin{definition}[Two endpoint absorption coefficients]\label{def:xi-endpoint-two-absorption-coefficients}
Let \(B\) be a finite Blaschke product with zero set \(\{a_k\}_{k=1}^{N}\subset\mathbb{D}\) (counted with multiplicity).
In addition to the endpoint \(-1\) absorption coefficient \(\Phi_{-1}(B)\) of Proposition \ref{prop:xi-endpoint-absorption-coefficient}, define the absorption coefficient at endpoint \(+1\)
\[
\Phi_{+1}(B):=\sum_{k=1}^{N}P_{a_k}(+1)=\sum_{k=1}^{N}\frac{1-|a_k|^2}{|1-a_k|^2}\ \in[0,\infty).
\]
\end{definition}

\begin{theorem}[Adams scale gate endpoint absorption parity law and linear amplification]\label{thm:xi-endpoint-absorption-adams-rescaling}
Let \(B\) be a finite Blaschke product, and let Adams scale gate \(\pi_m(w):=w^m\) (\(m\in\NN\)).
Denote the power pullback \(B^{(m)}:=B\circ\pi_m\), i.e., \(B^{(m)}(w)=B(w^m)\).
Then for any \(m\ge 1\),
\[
\boxed{\;
\Phi_{-1}\bigl(B^{(m)}\bigr)=
\begin{cases}
m\,\Phi_{-1}(B),& m\ \text{odd},\\[2pt]
m\,\Phi_{+1}(B),& m\ \text{even},
\end{cases}
\qquad
\Phi_{+1}\bigl(B^{(m)}\bigr)=m\,\Phi_{+1}(B).\;}
\]
\end{theorem}

\begin{proof}
First deduce the radial absorption law at endpoint \(+1\) from Proposition \ref{prop:xi-endpoint-absorption-coefficient}: let \(\widetilde B(w):=B(-w)\), then \(\widetilde B\) is still a finite Blaschke product and \(\widetilde B(-r)=B(r)\).
Moreover
\[
\Phi_{-1}(\widetilde B)
:=\sum_{k=1}^{N}\frac{1-|(-a_k)|^2}{|1+(-a_k)|^2}
=\sum_{k=1}^{N}\frac{1-|a_k|^2}{|1-a_k|^2}
=\Phi_{+1}(B),
\]
so by Proposition \ref{prop:xi-endpoint-absorption-coefficient},
\[
\log|B(r)|
=\log|\widetilde B(-r)|
=-(1-r)\,\Phi_{-1}(\widetilde B)+o(1-r)
=-(1-r)\,\Phi_{+1}(B)+o(1-r)
\qquad(r\uparrow 1).
\]
\par
For \(\Phi_{-1}(B^{(m)})\), take \(w=-r\uparrow-1\), then
\(
B^{(m)}(-r)=B((-r)^m)
\).
If \(m\) is odd, then \((-r)^m=-r^m\uparrow-1\); by Proposition \ref{prop:xi-endpoint-absorption-coefficient} (applied to \(r^m\uparrow 1\)),
\[
\log|B^{(m)}(-r)|=\log|B(-r^m)|=-(1-r^m)\,\Phi_{-1}(B)+o(1-r^m).
\]
If \(m\) is even, then \((-r)^m=r^m\uparrow+1\); by the \(+1\) radial absorption law from the previous step,
\[
\log|B^{(m)}(-r)|=\log|B(r^m)|=-(1-r^m)\,\Phi_{+1}(B)+o(1-r^m).
\]
Since \(1-r^m=m(1-r)+o(1-r)\) (\(r\uparrow 1\)), and \(B^{(m)}\) is still a finite Blaschke product to which Proposition \ref{prop:xi-endpoint-absorption-coefficient} applies, we have
\(
\log|B^{(m)}(-r)|=-(1-r)\,\Phi_{-1}(B^{(m)})+o(1-r)
\);
comparing first-order coefficients yields the stated parity law.
\par
For \(\Phi_{+1}(B^{(m)})\), take \(w=r\uparrow+1\), then \(B^{(m)}(r)=B(r^m)\) and \(r^m\uparrow 1\).
By the \(+1\) radial absorption law applied to \(B\) and \(B^{(m)}\) respectively,
\[
\log|B^{(m)}(r)|
=\log|B(r^m)|
=-(1-r^m)\,\Phi_{+1}(B)+o(1-r^m)
=-(1-r)\,m\,\Phi_{+1}(B)+o(1-r)
\]
and
\(
\log|B^{(m)}(r)|=-(1-r)\,\Phi_{+1}(B^{(m)})+o(1-r)
\).
Comparing first-order coefficients gives \(\Phi_{+1}(B^{(m)})=m\Phi_{+1}(B)\).\qed
\end{proof}

\begin{remark}[Parity and degeneration/merging of endpoint channels]\label{rem:xi-endpoint-absorption-adams-parity}
The power gate \(\pi_m(w)=w^m\) satisfies \(\pi_m(+1)=+1\) and \(\pi_m(-1)=(-1)^m\) on endpoints:
for odd \(m\) endpoint \(-1\) remains fixed, hence the \(\Phi_{-1}\) channel is linearly amplified without changing basepoint;
for even \(m\) endpoint \(-1\) is sent to \(+1\), so \(\Phi_{-1}(B^{(m)})\) reads the absorption coefficient \(\Phi_{+1}(B)\) of \(B\) at \(+1\).
In particular, for even \(m\) the two endpoint channels merge to the same indicator on \(B^{(m)}\):
\[
\Phi_{-1}(B^{(m)})=\Phi_{+1}(B^{(m)})=m\,\Phi_{+1}(B),
\]
hence endpoint audit degenerates to one dimension at this scale.
\end{remark}

\begin{remark}[Auditable terminal discriminant form: endpoint one-dimensional index and multiscale audit chain]\label{rem:xi-endpoint-absorption-terminal-index}
Under Hardy/Smirnov factorization and the calibration "offset only carried by Blaschke data," off-slice offset can be compressed into the endpoint one-dimensional index \(\Phi_{-1}(B)\) (together with endpoint \(+1\) index \(\Phi_{+1}(B)\), Definition \ref{def:xi-endpoint-two-absorption-coefficients}):
on one hand it equals the single-point radial decay rate limit \(\lim_{r\uparrow 1}\frac{-\log|B(-r)|}{1-r}\) (Corollary \ref{cor:xi-endpoint-absorption-additive-tomography}), on the other hand satisfies \(\Phi_{-1}(B)=0\iff B\equiv 1\) (Proposition \ref{prop:xi-endpoint-absorption-coefficient}), thus providing an endpoint discrimination interface for "no off-slice offset" independent of radial gap resolution.
Furthermore, Adams scale gate \(\pi_m\) gives deterministic linear rescaling of endpoint indices (Theorem \ref{thm:xi-endpoint-absorption-adams-rescaling}):
odd \(m\) preserves the endpoint \(-1\) channel, hence the family \(\{\Phi_{-1}(B^{(m)})\}_{m\ \mathrm{odd}}\) forms a lossless multiscale audit chain; it performs deterministic amplification of "total offset" as an additive conservation quantity with integer-coefficient scale parameters, without introducing new continuous degrees of freedom.
Even \(m\) triggers degeneration and merging of endpoint channels (Remark \ref{rem:xi-endpoint-absorption-adams-parity}), causing this chain to lose distinction between \(-1\) and \(+1\) at that scale.
\end{remark}

\begin{theorem}[Multiscale consistency \texorpdfstring{$\Longrightarrow$}{=>} fixed slice support: sufficient condition for odd Adams gate endpoint budget uniform upper bound]\label{thm:xi-multiscale-consistency-implies-rh}
Under Hardy/Smirnov factorization calibration, suppose the disk completion object \(F_\zeta\) has inner factor as pure Blaschke object \(B\) (no singular inner factor).
If there exists constant \(M<\infty\) such that for all odd \(m\ge 1\) Adams scale gates \(\pi_m(w)=w^m\) there is uniform upper bound
\[
\Phi_{-1}\bigl(B\circ \pi_m\bigr)\le M,
\]
then \(\Phi_{-1}(B)=0\), hence \(B\equiv 1\).
Therefore \(F_\zeta\) has no zeros in the open unit disk \(\mathbb{D}\), and RH holds (Corollary \ref{cor:app-endpoint-phase-certificate} and \S\ref{subsec:unit-disk-inner-outer}).
\end{theorem}

\begin{proof}
By Theorem \ref{thm:xi-endpoint-absorption-adams-rescaling}, for odd \(m\),
\(
\Phi_{-1}(B\circ\pi_m)=m\,\Phi_{-1}(B)
\).
If this is uniformly bounded for all odd \(m\), then letting \(m\to\infty\) gives \(\Phi_{-1}(B)=0\).
By the last sentence of Proposition \ref{prop:xi-endpoint-absorption-coefficient}, \(B\equiv 1\).
Under the no singular inner factor assumption, \(B\equiv 1\) is equivalent to no point-like zero data in the disk, hence \(F_\zeta\) has no zeros in \(\mathbb{D}\); by Corollary \ref{cor:app-endpoint-phase-certificate} this returns to RH.\qed
\end{proof}

\paragraph{Horizon renormalization: strict scaling law of endpoint flux}
The endpoint absorption coefficient \(\Phi_{-1}(B)\) is a single-point compression of in-disk point-like defects; the paired scale semantics is given by the "dilation semigroup."
Appendix \eqref{eq:scale-psi-m} introduces the Möbius semigroup \(\{\psi_m\}_{m>0}\subset\mathrm{Aut}(\mathbb{D})\) fixing \(\{\pm 1\}\), whose endpoint multiplier satisfies \(\psi_m'(-1)=m^{-1}\) (Proposition \ref{prop:scale-linearization-identity-multipliers}).
The following theorem shows that endpoint Poisson weight (i.e., single-atom contribution of flux) strictly amplifies by factor \(m\) under this scale chain.

\begin{theorem}[Scale covariance of endpoint Poisson weight]\label{thm:xi-endpoint-poisson-weight-scaling}
Let \(m>1\) and \(\psi_m\) as in \eqref{eq:scale-psi-m}.
For the endpoint Poisson weight of any \(a\in\mathbb{D}\)
\(
P_a(-1):=\frac{1-|a|^2}{|1+a|^2}
\)
there holds the strict identity
\[
\boxed{\;P_{\psi_m(a)}(-1)=m\,P_a(-1)\;}
\]
Therefore if \(B\) is a finite Blaschke product with zero set \(\{a_k\}_{k=1}^{N}\subset\mathbb{D}\) (counted with multiplicity), and letting
\[
a_k^{(n)}:=\psi_{m^n}(a_k),
\]
then letting \(B^{(n)}\) be the Blaschke product with zero set \(\{a_k^{(n)}\}\) (multiplicity preserved), the endpoint hair flux satisfies
\[
\boxed{\;\Phi_{-1}(B^{(n)})=m^n\,\Phi_{-1}(B)\;}
\]
where \(\Phi_{-1}\) is defined by Proposition \ref{prop:xi-endpoint-absorption-coefficient}.
\end{theorem}

\begin{proof}
Direct calculation from \eqref{eq:scale-psi-m} (also see proof of Proposition \ref{prop:scale-linearization-identity-multipliers}) yields the identity
\[
\frac{1-\psi_m(a)}{1+\psi_m(a)}=m\,\frac{1-a}{1+a}.
\]
On the other hand, for any \(z\in\mathbb{D}\),
\[
\Re\!\left(\frac{1-z}{1+z}\right)=\frac{1-|z|^2}{|1+z|^2}=P_z(-1),
\]
so taking real parts immediately gives \(P_{\psi_m(a)}(-1)=m\,P_a(-1)\).
The last formula follows by term-by-term summation and the semigroup property \(\psi_{m^n}=\psi_m^{\circ n}\).\qed
\end{proof}

\begin{corollary}[Cross-scale closure \texorpdfstring{$\Longrightarrow$}{=>} no-hair fixed point]\label{cor:xi-nohair-fixedpoint-flux-closure}
Let \(m>1\) and take the scale pushforward family \(\{B^{(n)}\}_{n\ge 0}\) from the above theorem.
If there exists constant \(M<\infty\) such that
\[
\sup_{n\ge 0}\Phi_{-1}(B^{(n)})\le M,
\]
then \(\Phi_{-1}(B)=0\), hence \(B\equiv 1\).
\end{corollary}

\begin{proof}
By Theorem \ref{thm:xi-endpoint-poisson-weight-scaling}, \(\Phi_{-1}(B^{(n)})=m^n\Phi_{-1}(B)\).
When \(m>1\), if the left side is uniformly bounded in \(n\), then necessarily \(\Phi_{-1}(B)=0\).
By the last sentence of Proposition \ref{prop:xi-endpoint-absorption-coefficient}, \(B\equiv 1\).\qed
\end{proof}

\begin{corollary}[Extreme resonance criterion via endpoint flux (no singular inner factor calibration)]\label{cor:xi-extreme-resonance-criterion}
Under Hardy/Smirnov condition calibration, write the disk completion object as factorization \(F_{\zeta}=B\,S\,O\) (Appendix \S\ref{subsec:unit-disk-inner-outer}).
Further assume no singular inner factor \(S\equiv 1\), and \(B\) is a finite Blaschke product (zero set counted with multiplicity as \(\{a_k\}\)).
Then the following propositions are equivalent:
\begin{enumerate}
  \item \(B\equiv 1\) (no in-disk point-like zero data);
  \item \(\Phi_{-1}(B)=0\);
  \item Endpoint phase current density \(j_{\mathrm{in}}(\pi)=0\), where \(j_{\mathrm{in}}(\pi):=\sum_k P_{a_k}(-1)\) (Theorem \ref{thm:app-inner-phase-current-poisson}).
\end{enumerate}
And if there exists auditable upper bound \(j_{\mathrm{in}}(\pi)\le M<\infty\), then the total off-critical-line offset satisfies linear constraint
\[
\sum_k\delta_k\le M.
\]
\end{corollary}

\begin{proof}
Under \(S\equiv 1\) and \(B\) being a finite Blaschke product, by Theorem \ref{thm:app-inner-phase-current-poisson},
\(
j_{\mathrm{in}}(\pi)=\sum_k P_{a_k}(-1)=\Phi_{-1}(B)
\).
Each term \(P_{a_k}(-1)\) is strictly positive, so
\(
j_{\mathrm{in}}(\pi)=0
\iff
\Phi_{-1}(B)=0
\iff
\{a_k\}=\varnothing
\iff
B\equiv 1
\), giving equivalence.
The last formula follows directly from Corollary \ref{cor:app-endpoint-phase-current-offcritical} (or the upper bound paragraph of Corollary \ref{cor:app-endpoint-phase-certificate}): \(j_{\mathrm{in}}(\pi)\ge \sum_k\delta_k\), so if \(j_{\mathrm{in}}(\pi)\le M\) then \(\sum_k\delta_k\le M\).\qed
\end{proof}

\begin{remark}[Mathematical formulation of threshold carrying (extreme resonance)]\label{rem:xi-threshold-carrying-extreme-resonance}
In this paper's disk-fixed slice semantics, the existence of finite Blaschke data \(B\) is equivalent to the point spectrum residue axis of minimal realization: its minimal residue dimension is \(\deg B\) (Theorem \ref{thm:blaschke-point-spectrum-correspondence}).
Therefore once \(B\equiv 1\) (i.e., endpoint flux is zero), in-disk discrete point spectrum resonances are excluded;
if singularity is still allowed, it can only be carried in boundary form on the fixed slice \(|w|=1\) (e.g., via boundary singular inner factor \(S\) and boundary measure, see Theorem \ref{thm:blaschke-defect-zero-fixed-slice-support} and Appendix \S\ref{subsec:unit-disk-inner-outer}).
\par
In this sense, "extreme resonance" can be strictly understood as: singularity no longer appears as in-disk discrete modes, but only as threshold boundary support;
this formulation is compatible with the unitary slice closure semantics (refusing additional orthogonal residue axes) and strict Beck-Chevalley/strictification zero anomaly calibration (refer to the reduction chain of Theorem \ref{thm:dephys-offline-hair-curvature-reduction}).
\end{remark}

\begin{proposition}[Resonance quality factor and Toeplitz-PSD detection threshold]\label{prop:xi-extreme-resonance-toeplitz-threshold}
Following the horizon Carathéodory logarithmic derivative \(\mathscr{C}_{\zeta}\) and Toeplitz-PSD hierarchy (Appendix Theorem \ref{thm:app-horizon-toeplitz-lmi}).
If there exists off-critical-line zero \(\rho=\beta+\mathrm{i}\gamma\) (\(\beta<\tfrac12\), \(\delta:=\tfrac12-\beta>0\)), denote its disk image \(w_\rho=\mathcal{C}(\gamma+\mathrm{i}\delta)\in\mathbb{D}\), and define resonance quality factor
\[
Q_\rho:=\frac{1}{1-|w_\rho|^2}=\frac{\gamma^2+(1+\delta)^2}{4\delta}
\qquad\text{(Corollary \ref{cor:app-offcritical-horocycle-busemann})}.
\]
Then in-disk poles leave explicit leading-term imprints in the multiple moments \(\{c_n\}\) of \(\mathscr{C}_{\zeta}\) (Proposition \ref{prop:app-horizon-pole-imprint}), hence under the regime of "remainder has auditable upper bound," the finite-level verification of Toeplitz-PSD has an intrinsic detection threshold: if one wishes to exclude this pole by finite-level PSD certificate (i.e., force a finite-dimensional failure witness), the required moment order \(N\) must satisfy
\[
N\ \gtrsim\ Q_\rho\,\log\!\bigl(1+\delta Q_\rho\bigr)
\asymp \frac{\gamma^2}{\delta}\,\log(1+\gamma^2)
\]
(Appendix Theorem \ref{thm:app-horizon-toeplitz-detection-threshold}; see also Remark \ref{rem:app-horizon-toeplitz-witness-endpoint-focus} for endpoint focusing interpretation).
\end{proposition}

\begin{corollary}[Quadratic compression under height window: Toeplitz truncation order and \texorpdfstring{\(\varphi\)}{phi}-calibration law]\label{cor:xi-toeplitz-height-phi-calibration}
Following the notation of Proposition \ref{prop:xi-extreme-resonance-toeplitz-threshold}, let
\(
m_\star(\rho):=\lceil\log_{\varphi}Q_\rho\rceil
\)
be the \(\varphi\)-resolution threshold determined by horizon depth \(1-|w_\rho|^2\) (Theorem \ref{thm:xi-bulk-horizon-phi-area-law}).
Then for finite-level Toeplitz-PSD verification to produce a finite-dimensional failure witness for this pole, the required moment order satisfies
\[
N\ \gtrsim\ \varphi^{m_\star(\rho)}\,\log\!\bigl(1+\delta\,\varphi^{m_\star(\rho)}\bigr)
\qquad(\delta:=\tfrac12-\beta>0).
\]
Hence for fixed moment order \(N\) truncation verification, poles that can be forcibly detected can only satisfy
\(
m_\star(\rho)\ \lesssim\ \log_{\varphi}N
\)
(at most a logarithmic factor difference).
\par
More concretely, given offset lower bound \(\delta\ge\delta_0>0\) and height window \(|\gamma|\le T\), denote the worst quality factor in this window
\[
Q_{\max}(T,\delta_0):=\frac{T^2+(1+\delta_0)^2}{4\delta_0},
\qquad
m_{\max}(T,\delta_0):=\left\lceil\log_{\varphi}Q_{\max}(T,\delta_0)\right\rceil.
\]
Then any finite-level Toeplitz-PSD certificate wishing to \emph{uniformly} exclude potential off-slice poles within this height window must satisfy
\[
N\ \gtrsim\ Q_{\max}(T,\delta_0)\,\log\!\bigl(1+\delta_0 Q_{\max}(T,\delta_0)\bigr)
\asymp \frac{T^2}{\delta_0}\,\log(1+T^2),
\]
and the corresponding \(\varphi\)-resolution threshold satisfies
\[
m_{\max}(T,\delta_0)=2\log_{\varphi}T+O(1)\qquad(T\to\infty).
\]
\end{corollary}
\begin{proof}
By definition \(m_\star(\rho)=\lceil\log_{\varphi}Q_\rho\rceil\), we have \(Q_\rho=\Theta(\varphi^{m_\star(\rho)})\); substituting into Proposition \ref{prop:xi-extreme-resonance-toeplitz-threshold} yields the first formula.
The reverse estimate follows directly from the same formula: if \(m_\star(\rho)\) is much larger than \(\log_\varphi N\), the right-hand leading term strictly exceeds \(N\), and finite-level verification cannot forcibly produce a failure witness.
\par
Finally, for \(\delta\in(0,1)\),
\(
Q_\rho=(\gamma^2+1)/(4\delta)+\tfrac12+\delta/4
\)
increases monotonically with \(|\gamma|\) and decreases monotonically with \(\delta\), so under constraints \(|\gamma|\le T\), \(\delta\ge\delta_0\) we have
\(
Q_\rho\le Q_{\max}(T,\delta_0)
\)
with worst case achieved at \((|\gamma|,\delta)=(T,\delta_0)\).
Substituting \(Q_{\max}\) into Proposition \ref{prop:xi-extreme-resonance-toeplitz-threshold} gives the uniform lower bound; and \(m_{\max}=\lceil\log_\varphi Q_{\max}\rceil\) gives the last formula.\qed
\end{proof}

Finally it must be emphasized that "zero leakage" should be understood as a verifiable regime rather than single-point coincidence.
On the disk completion \(F_\zeta\), RH is equivalent to defect identically zero \(D_{\zeta}(\varrho)\equiv 0\) (Corollary \ref{cor:app-rh-iff-jensen-defect-zero}), also equivalent to the countable constraint \(D_{\zeta}(\varrho_n)=0\) for any approaching sequence \(\varrho_n\uparrow 1\) (Theorem \ref{thm:app-jensen-countable-criterion}).
Under this "zero leakage regime," the Jensen-Blaschke mechanism excludes all in-disk point-like zeros; combined with the fixed slice \(|u|=1\iff \Re(s)=\tfrac12\) of completion self-dual parameter \(u=u(s)\), we obtain: point-like zeros can only be supported on the fixed slice, hence can only be supported on \(\Re(s)=\tfrac12\).
Paired with this, the role of Ramanujan-type symmetry in this framework is not "axis exchange," but forcing phase-closing through local spectral bounds, thereby channel-wise eliminating amplitude/radial leakage (\S\ref{subsubsec:ramanujan-phase-lift-symmetry}; Theorem \ref{thm:ramanujan-phase-lift-principle}).
\par
Therefore in this paper's system, the unitary slice is the minimal section of continuous visible parameters: radial information off this section enters at the endpoint layer with quadratic compression law (Theorem \ref{thm:xi-offcritical-quadratic-radial-compression}), its visibilization requires irreducible endpoint resolution/bit budget (Theorem \ref{thm:xi-offcritical-endpoint-resolution-lower-bound}), and can be certificated on finite height bands by single-radius Jensen defect inequality certificate (Theorem \ref{thm:xi-jensen-single-radius-band-exclusion}).
Comoving localization can develop the offset at linear scale, but its cost is introducing an addressable prefix that must be explicitly given (Remark \ref{rem:xi-comoving-horizon-localization-address-null}); under fixed visible interface if one refuses to externalize this prefix, off-slice offset is still forced through the endpoint second-order radial channel and triggers corresponding capacity threshold.
The above budget axes are mutually orthogonal and non-substitutable in protocol semantics (Theorem \ref{thm:xi-null-complete-trichotomy-offline}): raising budget on one axis cannot cancel threshold violation on another; when refusing to externalize the mode axis, "off-slice assertions" cannot form verifiable addresses at the type level and stably degrade to \(\mathsf{NULL}\).

\endinput
